%--------------------------------------------------------------------------------
% Szablon pracy dyplomowej inżynierskiej i magisterskiej
% Opracowanie: Jakub Stankiewicz
% Politechnika Warszawska, 2024-2025
% 
% Ten szablon jest udostępniony na licencji Creative 
% Commons Attribution-ShareAlike 4.0 International (CC BY-SA 4.0).
% Masz prawo do:
% - Udostępniania: kopiowania i rozpowszechniania materiału na dowolnym nośniku
%   i w dowolnym formacie.
% - Adaptacji: remiksowania, przekształcania i tworzenia nowych materiałów na 
%   podstawie materiału, również do celów komercyjnych.
%
% Na następujących warunkach:
% - Uznanie autorstwa: musisz odpowiednio oznaczyć autorstwo, podać link do licencji
%   oraz wskazać, czy wprowadziłeś zmiany. Możesz to zrobić w dowolny sposób,
%   ale nie w sposób sugerujący, że licencjodawca popiera ciebie lub twoje użycie.
% - Na tych samych warunkach: jeśli remiksujesz, przekształcasz lub tworzysz nowe 
%   materiały na podstawie oryginału, musisz rozpowszechniać swoje dzieło na tej 
%   samej licencji, co oryginał.
% 
% Szczegóły licencji: https://creativecommons.org/licenses/by-sa/4.0/
%
% Zrzeczenie się odpowiedzialności:
% Szablon jest udostępniony "tak jak jest" bez żadnej gwarancji lub zapewnienia
% przydatności do określonego celu. Autor nie ponosi odpowiedzialności za 
% jakiekolwiek problemy wynikające z jego użycia.
%--------------------------------------------------------------------------------
%  - Bibliografia: Kamil Zdanowicz, Krzysztof Dobrzyniecki  
%  - Strona tytułowa: Rafał Łysoń
%  - Treść szablonu: Jakub Stankiewicz
%--------------------------------------------------------------------------------
%  Sposób kompilacji:`'
%  pdflatex -> biber -> pdflatex x2
%--------------------------------------------------------------------------------
%  MiKTeX-pdfTeX 4.19 (MiKTeX 24.4)
%  Wersja biber: 2.20
%--------------------------------------------------------------------------------
% Czcionki w LaTeX:
% +---------------+--------------------+--------------------+--------------------+
% | Nazwa         | Rozmiar dla 10pt  | Rozmiar dla 11pt   | Rozmiar dla 12pt   |
% +---------------+--------------------+--------------------+--------------------+
% | \tiny         | 5pt               | 6pt                | 6pt                |
% | \scriptsize   | 7pt               | 8pt                | 8pt                |
% | \footnotesize | 8pt               | 9pt                | 10pt               |
% | \small        | 9pt               | 10pt               | 10.95pt            |
% | \normalsize   | 10pt              | 10.95pt            | 12pt               |
% | \large        | 12pt              | 12pt               | 14.4pt             |
% | \Large        | 14.4pt            | 14.4pt             | 17.28pt            |
% | \LARGE        | 17.28pt           | 17.28pt            | 20.74pt            |
% | \huge         | 20.74pt           | 20.74pt            | 24.88pt            |
% | \Huge         | 24.88pt           | 24.88pt            | 29.86pt            |
% +---------------+--------------------+--------------------+--------------------+

% !TEX root = main.tex
\documentclass[a4paper, final , 11pt, twoside, openany]{book} 

%-----------------------------------Bibliografia---------------------------------
\RequirePackage[backend=biber,datamodel=thesis,style=numeric,giveninits=true]{biblatex}
\usepackage{xpatch}

% Przecinki zamiast kropek pomiędzy polami
\DeclareDelimFormat{newunitpunct}{\addcomma\space}
% Brak kropki na końcu wpisu
\DeclareDelimFormat{finentrypunct}{}
% Lokacje oddzielamy slashem
\DeclareDelimFormat{multilistdelim}{\addslash}
\DeclareDelimFormat{finallistdelim}{\addslash}
\DeclareDelimFormat{bibdaterangesep}{\textendash}

% Kolejność nazwisko-imię lub na odwrót 
\DeclareNameAlias{author}{family-given}
\DeclareNameAlias{editor}{given-family}
\DeclareNameAlias{translator}{given-family}
\DeclareNameAlias{default}{given-family}
\DeclareNameAlias{supervisor}{given-family}

% Tytuł jest bez formatowania, chyba że jest
\DeclareFieldFormat*{title}{#1}
\DeclareFieldFormat[book]{title}{\mkbibemph{#1}}

% strony, dla książek w nawiasie
\DeclareFieldFormat{pages}{ss.\nobreakspace#1}
\DeclareFieldFormat[book]{pages}{\mkbibparens{ss.\nobreakspace#1}}
\DeclareFieldFormat[inbook]{pages}{\mkbibparens{ss.\nobreakspace#1}}

% Trochę zmieniamy teksty wbudowane
\NewBibliographyString{engthesis}
\DefineBibliographyStrings{polish}{%
    phdthesis   =   {rozprawa doktorska},
    engthesis   =   {praca inżynierska},
}

\xpatchbibmacro{byeditor+others}{\usebibmacro{bytranslator+others}}
{\usebibmacro{bytranslator}}{}{\ddt}

% tłumacz w nawiasie
\renewbibmacro*{bytranslator}{%
    \ifnameundef{translator}
    {}
    {\addcomma\space\mkbibparens{tłum.
        \setunit{\addspace}%
        \printnames[bytranslator]{translator}}
    }
}

% wydanie
\DeclareFieldFormat{edition}{\mkbibparens{wyd.\nobreakspace#1}}
\renewbibmacro*{url+urldate}{%
    \iffieldundef{url}{}
    {
        \mkbibacro{URL}\nopunct
        \iffieldundef{urlyear}{}
        {
            \setunit*{\addspace}%
            \usebibmacro{urldate}
        }\addcolon\addspace
        \usebibmacro{url}%
    }
}
\DeclareFieldFormat{url}{\url{#1}}
\DeclareFieldFormat{urldate}{\mkbibparens{dost. #1}}
\DeclareFieldFormat{indextitle}{#1,}
\DeclareFieldFormat{note}{\mkbibparens{#1}}

% Zmienia "w" na "W:"
\renewbibmacro*{in:}{%
  \bibcpstring{in}\addcolon%
  \printunit{\intitlepunct}
}

% Autorzy: max 7, ogranicza do 7. Redaktorzy: Max 3, ogranicza do 1
\xpatchbibmacro{author}{\printnames{author}}{\printnames[][-7]{author}}{}{\ddt}
\xpatchbibmacro{editor}{\printnames{editor}}{%
    \printnames[][-\ifnum\value{listtotal}>3 1 \else 3 \fi]{editor}}{}{\ddt}
\xpatchbibmacro{editor+others}{\printnames{editor}}{%
    \printnames[][-\ifnum\value{listtotal}>3 1 \else 3 \fi]{editor}}{}{\ddt}
\xpatchbibmacro{byauthor}{\printnames[byauthor]{author}}{\printnames[byauthor][-7]{author}}{}{\ddt}
\xpatchbibmacro{byeditor}{\printnames[byeditor]{editor}}{%
    \printnames[][-\ifnum\value{listtotal}>3 1 \else 3 \fi]{editor}}{}{\ddt}

% Jeśli jest wydanie, to po tytule nie dajemy przecinka
\xpatchbibmacro{maintitle+title}{%
    \usebibmacro{title}%
}{%
    \usebibmacro{title}%
    \iffieldundef{edition}{}{\nopunct}%
}{}{\ddt}

% Do 10 lokalizacji
\xpatchbibmacro{institution+location+date}{\printlist{location}}{\printlist[][-10]{location}}{}{\ddt}
\xpatchbibmacro{organization+location+date}{\printlist{location}}{\printlist[][-10]{location}}{}{\ddt}
\xpatchbibmacro{location+date}{\printlist{location}}{\printlist[][-10]{location}}{}{\ddt}
\renewbibmacro*{date}{\printdate\adddot\nopunct}
% Zmieniamy trochę format
\renewbibmacro*{publisher+location+date}{
    \printlist{publisher}%
    \newunit\newblock
    \printlist[][-10]{location}%
    \newunit\newblock
    \usebibmacro{date}\isdot
}

% Umieszcza "red." w nawiasie
\renewbibmacro*{editor+othersstrg}{\addspace \mkbibparens{\bibxstring{editor}}}
\renewbibmacro*{editorstrg}{\addspace \mkbibparens{\bibxstring{editor}}}

\newbibmacro*{publisher+location}{%
    \printlist{publisher}%
    \newunit\newblock
    \printlist[][-10]{location}%
}

% Strony przed datą
\xpatchbibdriver{book}{%
    \usebibmacro{publisher+location+date}%
    \newunit\newblock
    \usebibmacro{chapter+pages}%
    \newunit
    \printfield{pagetotal}%
}{%
    \usebibmacro{publisher+location}%
    \newunit\newblock
    \usebibmacro{chapter+pages}%
    \newunit
    \printfield{pagetotal}%
    \newunit\newblock
    \usebibmacro{date}%
}{}{\ddt}

% Przesuwa pole 'note' zaraz po języku tytułu (w praktyce po tytule) w @book
\xpatchbibdriver{book}{%
    \printfield{note}%
    \newunit\newblock
}{}{}{\ddt}
\xpatchbibdriver{book}{%
    \printlist{language}%
    \newunit\newblock
}{%
    \printlist{language}%
    \newunit\newblock
    \printfield{note}%
    \newunit\newblock
}{}{\ddt}

% Redakcja przed nazwą książki w @inbook
\xpatchbibdriver{inbook}{%
    \usebibmacro{bybookauthor}%
    \newunit\newblock
    \usebibmacro{maintitle+booktitle}%
    \newunit\newblock
    \usebibmacro{byeditor+others}%
}{%
    \newunit\newblock
    \usebibmacro{editor+others}%
    \newunit\newblock
    \usebibmacro{maintitle+booktitle}%
}{}{\ddt}

% Numery po tytule książki w @inbook
\DeclareFieldFormat[inbook]{booktitle}{\mkbibemph{#1}\nopunct}
\xpatchbibdriver{inbook}{%
    \usebibmacro{maintitle+booktitle}%
}{%
    \usebibmacro{maintitle+booktitle}%
    \newunit\newblock
    \printfield{pages}\clearfield{pages}%
}{}{\ddt}

\DeclareFieldFormat[article]{journaltitle}{\mkbibemph{#1},\nopunct}
\DeclareFieldFormat[article]{volume}{\mkbibemph{#1}\nopunct}
\DeclareFieldFormat[proceedings]{volume}{#1\nopunct}
\DeclareFieldFormat[inproceedings]{volume}{#1\nopunct}
\DeclareFieldFormat[article]{number}{(#1)\nopunct}
\xpatchbibdriver{article}{%
    \usebibmacro{in:}%
    \usebibmacro{journal+issuetitle}%
    \newunit
    \usebibmacro{byeditor+others}%
    \newunit
    \usebibmacro{note+pages}%
    \newunit\newblock
}{%
    \usebibmacro{journal+issuetitle}%
    \newunit
    \printfield{pages}%
    \newunit
    \usebibmacro{byeditor+others}%
    \newunit
    \usebibmacro{date}
    \newunit
    \printfield{note}%
    \newunit\newblock
}{}{\ddt}
\xpatchbibmacro{journal+issuetitle}{%
    \usebibmacro{issue+date}%
}{%
    \usebibmacro{issue},%
}{}{\ddt}

\DeclareFieldFormat[thesis]{type}{\ifbibstring{#1}{\bibcplstring{#1}}{#1}}
\DeclareNameWrapperFormat*[thesis]{supervisor}{promotor #1}
\newbibmacro*{institution+location}{
    \printlist{institution}, \printlist{location}
}
\xpatchbibdriver{thesis}{%
    \newunit
    \usebibmacro{institution+location+date}%
    \newunit\newblock
}{%
    \newunit\newblock
    \usebibmacro{institution+location}%
    \newunit\newblock
    \printnames{supervisor}%
    \newunit\newblock
    \usebibmacro{date}\isdot%
    \newunit\newblock
}{}{\ddt}

\DeclareFieldFormat[standard]{title}{\mkbibemph{#1},}

\DeclareFieldFormat[manual]{version}{wer. #1,}
\DeclareListWrapperFormat[manual]{organization}{#1,}
%\xpatchbibdriver{manual}{%
%    \printlist{organization}%
%    \newunit
%}{}{}{\ddt}
\xpatchbibdriver{manual}{%
    \usebibmacro{author/editor}%
}{%
    \usebibmacro{author/editor}%
    \newunit\newblock
    \ifnameundef{author}{\ifnameundef{editor}{%
        \printlist{organization}
        \clearlist{organization}
        \newunit
    }{}} {}%
}{}{\ddt}
\xpatchbibdriver{manual}{%
    \usebibmacro{publisher+location+date}%
}{%
    \usebibmacro{date}%
}{}{\ddt}

\DeclareListWrapperFormat[standard]{organization}{#1,}
\DeclareBibliographyDriver{standard}{%
  \usebibmacro{bibindex}%
  \usebibmacro{begentry}%
  \printfield{indextitle}%
  \newunit\newblock
  \setunit{\printdelim{nametitledelim}}\newblock
  \usebibmacro{title}%
  \newunit
  \printlist{language}%
  \newunit\newblock
  \printlist{organization}%
  \newunit\newblock
  \usebibmacro{byauthor}%
  \newunit\newblock
  \printfield{note}%
  \newunit\newblock
  \printfield{type}%
  \newunit
  \usebibmacro{institution+location+date}%
  \newunit\newblock
  \usebibmacro{chapter+pages}%
  \newunit
  \printfield{pagetotal}%
  \newunit\newblock
  \iftoggle{bbx:isbn}
    {\printfield{isbn}}
    {}%
  \newunit\newblock
  \usebibmacro{doi+eprint+url}%
  \newunit\newblock
  \usebibmacro{addendum+pubstate}%
  \setunit{\bibpagerefpunct}\newblock
  \usebibmacro{pageref}%
  \newunit\newblock
  \iftoggle{bbx:related}
    {\usebibmacro{related:init}%
     \usebibmacro{related}}
    {}%
  \usebibmacro{finentry}}

\DeclareBibliographyDriver{legislation}{%
  \usebibmacro{bibindex}%
  \usebibmacro{begentry}%
  \setunit{\printdelim{nametitledelim}}\newblock
  \usebibmacro{title}%
  \newunit
  \printlist{language}%
  \newunit\newblock
  \printfield{booktitle}%
  \newunit\newblock
  \printfield{number}%
  \newunit\newblock
  \printfield{position}%
  \newunit\newblock
  \printfield{note}%
  \newunit\newblock
  \usebibmacro{finentry}}
\DeclareFieldFormat[legislation]{booktitle}{#1}
\DeclareFieldFormat[legislation]{number}{nr #1}
\DeclareFieldFormat[legislation]{position}{poz. #1}
\DeclareFieldFormat[legislation]{note}{#1}

\DeclareFieldFormat[online]{series}{#1,}
\DeclareFieldFormat[online]{volume}{nr #1}
\xpatchbibdriver{online}{%
    \printfield{version}%
    \newunit
}{%
    \printfield{series}%
    \newunit\newblock
    \printfield{volume}%
    \newunit\newblock
    \printfield{version}%
    \newunit
}{}{\ddt}

\DeclareFieldFormat[inproceedings]{booktitle}{\mkbibemph{#1},}
\DeclareFieldFormat[inproceedings]{maintitle}{\mkbibemph{#1},}
\DeclareNameWrapperFormat[inproceedings]{editor}{#1}
\xpatchbibdriver{inproceedings}{%
    \usebibmacro{chapter+pages}%
    \newunit\newblock
}{}{}{\ddt}
\xpatchbibdriver{inproceedings}{%
    \usebibmacro{maintitle+booktitle}%
    \newunit\newblock
    \usebibmacro{event+venue+date}%
    \newunit\newblock
}{%
    \usebibmacro{editor+others}%
    \newunit\newblock
    \usebibmacro{maintitle+booktitle}%
    \newunit\newblock
    \usebibmacro{event+venue+date}%
    \newunit\newblock
    \usebibmacro{chapter+pages}%
    \newunit\newblock
}{}{\ddt}
\xpatchbibmacro{event+venue+date}{%
    \printtext[parens]{%
        \printfield{venue}%
        \setunit*{\addcomma\space}%
    \printeventdate}%
}{%
    \printfield{venue}%
    \setunit*{\addcomma\space}%
    \printeventdate%
}{}{\ddt}

\addbibresource{bibliografia.bib}   % Plik z bibliografią

%-------------------------------------Pakiety-----------------------------------
\usepackage[
    font=small,
    labelfont=bf,
    labelsep=period,justification=justified,
    format=hang
]{caption}

\usepackage[
    a4paper,
    left=30mm,
    right=20mm,
    top=25mm,
    bottom=25mm
]{geometry}

\usepackage[
    protrusion=true,
    expansion=true,
    letterspace=400
]{microtype}

\usepackage[
    bookmarksnumbered=true,
    bookmarksopenlevel=3
]{hyperref}

\usepackage[normalem]{ulem} % do podkreślania tekstu - użyte w tabeli
\usepackage[hang]{footmisc}
\usepackage{listings}      % Code listings
\usepackage{fourier}       % Czcionka Adobe Utopia
\usepackage{amsmath}       % Basic mathematical typesetting
\usepackage{amssymb}       % Advanced math symbols
\usepackage{amsthm}        % Theorems typesetting
\usepackage{array}         % Advanced table column formats
\usepackage{enumitem}      % Itemize/enumerate
\usepackage{fancyhdr}      % Custom header/footer styles
\usepackage{graphicx}      % Enhanced images support
\usepackage{lipsum}        % Lorem ipsum
\usepackage{longtable}     % Multi-page tables
\usepackage{multirow}      % Advanced table cells
\usepackage{setspace}      % Set space between lines
\usepackage{scrextend}     % Allows \addmargin environment
\usepackage{tablefootnote} % Table footnotes
\usepackage{tocloft}       % Custom ToC/LoF/LoT
\usepackage{url}           % URL-sensitive line breaks
\usepackage{xspace}        % Remove duplicated spaces
\usepackage{layout}
\usepackage{booktabs}
\usepackage{etoolbox}
\usepackage{subcaption}     % do podwójnych tabel
\usepackage{float}
\usepackage{titlesec}


\usepackage{acro}
\usepackage[polish]{babel}
\usepackage{csquotes}


%--------------------------- Redefinicja nazw -----------------------------------

\renewcommand{\listfigurename}{Spis rysunków}
\renewcommand{\listtablename}{Spis tabel}
\renewcommand{\tablename}{Tabela}
\renewcommand{\contentsname}{Spis treści}
\renewcommand{\lstlistingname}{Kod źródłowy}
\acsetup{list/name=Akronimy}

%--------------------------- Ciągła numeracja -----------------------------------

\renewcommand{\thefigure}{\arabic{figure}} %numeracja ciągła
\renewcommand{\thetable}{\arabic{table}}
\renewcommand{\theequation}{\arabic{equation}}
\lstset{numberbychapter=false}


%--------------------- Ustawienia obiektów typu FLOAT ----------------------------
\setlength{\footnotemargin}{4mm} 
\setlength{\parindent}{5mm}      % Wcięcie akapitowe
\setlength{\textfloatsep}{10pt}  % Odstęp między floatem na górze lub dole strony a tekstem.
\setlength{\floatsep}{10pt}      % Odstęp między dwoma floatami w pionie.
\setlength{\intextsep}{10pt}     % Odstęp między floatem wstawionym w tekście (np. [h]) a tekstem.

%--------------------------- Inne ustawienia -----------------------------------
\setcounter{tocdepth}{3}
\setcounter{secnumdepth}{3}


%-------------------------------- Spis treści -----------------------------------
\setlength{\cftpartnumwidth}{3em}           % Szerokość miejsca na numer części w spisie treści
\renewcommand{\cftchapaftersnum}{.}         % Kropka po numerze rozdziału
\setlength{\cftchapnumwidth}{2.5em}         % Szerokość miejsca na numer rozdziału w spisie treści
\renewcommand{\cftchapleader}{\cftdotfill{\cftdotsep}} % kropki w spisie treści dla chapter
\renewcommand{\cftsecaftersnum}{.}          % Kropka po numerze sekcji
\setlength{\cftsecnumwidth}{3em}            % Szerokość miejsca na numer sekcji w spisie treści
\renewcommand{\cftsubsecaftersnum}{.}       % Kropka po numerze podsekcji
\setlength{\cftsubsecnumwidth}{2.5em}       % Szerokość miejsca na numer podsekcji w spisie treści
\renewcommand{\cftsubsubsecaftersnum}{.}    % Kropka po numerze podpodsekcji
\setlength{\cftsubsubsecnumwidth}{3.5em}    % Szerokość miejsca na numer podpodsekcji w spisie treści
\renewcommand{\cfttabaftersnum}{.}          % Kropka po numerze tabeli
\setlength{\cfttabnumwidth}{1.5em}          % Szerokość miejsca na numer tabeli w spisie treści
\renewcommand{\cftfigaftersnum}{.}          % Kropka po numerze ilustracji
\setlength{\cftfignumwidth}{1.5em}          % Szerokość miejsca na numer ilustracji w spisie treści


%-------------------------------- Redefinicja części -------------------------------
% Umieszcza part na odrębnej stronie z numerem części ponad nazwą części
% Dodaje słowo "część" do spisu treści 
\renewcommand{\part}[1]{
    \cleardoublepage
    \thispagestyle{empty} % Usuwa numer strony
    \refstepcounter{part} % Zwiększa licznik części
    \addcontentsline{toc}{part}{Część \thepart. #1}
    \vspace*{\fill} % Wyrównanie do środka strony (góra-dół)
    {\centering
        \normalfont\huge\bfseries Część \thepart. \\ % Nagłówek "Część <numer>"
        \vspace{0.5cm} % Odstęp pod nagłówkiem
        \Huge #1\par % Tytuł części
    }%
    \vspace*{\fill} % Wyrównanie dolnej części strony
    \newpage % Przechodzi na nową stronę
}

% ---------------------- Formatowanie rozdziałów ----------------------

% Rozdział numerowany
\titleformat{name=\chapter} % Format dla numerowanych rozdziałów
  {\normalfont\huge\bfseries} % Styl tytułu
  {\thechapter.\hspace{1em}} % Numer rozdziału
  {0pt} % Odstęp
  {\huge} % Styl tekstu tytułu

% Rozdział nienumerowany
\titleformat{name=\chapter,numberless} % Format dla nienumerowanych rozdziałów
  {\normalfont\huge\bfseries} % Styl tytułu
  {} % Brak numeru
  {0pt} % Odstęp
  {\huge} % Styl tekstu tytułu

% Ustawienia odstępów dla rozdziałów
\titlespacing*{\chapter}{0pt}{-20pt}{40pt}

% ---------------------- Formatowanie sekcji ----------------------
% Sekcja numerowana
\titleformat{name=\section} % Format dla numerowanych sekcji
  {\Large\bfseries} % Styl tytułu
  {\thesection.\hspace{1em}} % Numer sekcji
  {0.3cm} % Odstęp
  {}

% Sekcja nienumerowana
\titleformat{name=\section,numberless} % Format dla nienumerowanych sekcji
  {\Large\bfseries} % Styl tytułu
  {} % Brak numeru
  {0.3cm} % Odstęp
  {}

% Ustawienia odstępów dla sekcji
\titlespacing{\section}{0pt}{0.5cm}{0.5cm}

% ---------------------- Formatowanie podsekcji ----------------------
% Podsekcja numerowana
\titleformat{name=\subsection} % Format dla numerowanych podsekcji
  {\large\bfseries} % Styl tytułu
  {\thesubsection.\hspace{1em}} % Numer podsekcji
  {0.1cm} % Odstęp
  {}

% Podsekcja nienumerowana
\titleformat{name=\subsection,numberless} % Format dla nienumerowanych podsekcji
  {\large\bfseries} % Styl tytułu
  {} % Brak numeru
  {0.1cm} % Odstęp
  {}

% Ustawienia odstępów dla podsekcji
\titlespacing{\subsection}{0pt}{0.4cm}{0.4cm}

% ---------------------- Formatowanie podpodsekcji ----------------------
% Podpodsekcja numerowana
\titleformat{name=\subsubsection} % Format dla numerowanych podpodsekcji
  {\normalsize\bfseries} % Styl tytułu
  {\thesubsubsection.\hspace{1em}} % Numer podpodsekcji
  {0.1cm} % Odstęp
  {}

% Podpodsekcja nienumerowana
\titleformat{name=\subsubsection,numberless} % Format dla nienumerowanych podpodsekcji
  {\normalsize\bfseries} % Styl tytułu
  {} % Brak numeru
  {0.1cm} % Odstęp
  {}

% Ustawienia odstępów dla podpodsekcji
\titlespacing{\subsubsection}{0pt}{0.3cm}{0.3cm}


%--------------------------------- Formatowanie list ---------------------------------

% Odstępy między tekstem a listami, marginesami oraz czcionki
\setlist[itemize,1]{topsep=3pt,itemsep=1.5pt, parsep=2pt,label=\large$\bullet$, leftmargin=24pt}
\setlist[itemize,2]{topsep=3pt,itemsep=1pt, parsep=2pt,leftmargin=18pt, label=--}
\setlist[enumerate,1]{label=\arabic*.,topsep=3pt,itemsep=1.5pt, parsep=2pt,leftmargin=24pt}
\setlist[enumerate,2]{label=\alph*),topsep=3pt,itemsep=1pt, parsep=2pt,leftmargin=18pt}
\setlist[description,1]{labelindent=12pt, topsep=3pt,itemsep=1.5pt, parsep=2pt}
 % dla standardowych, krótkich list nie zmieniamy czcionki. Zmieniamy ręcznie dla list długich
\setlist[itemize]{before=\normalsize}     
\setlist[enumerate]{before=\normalsize}   
\setlist[description]{before=\normalsize}

\renewlist{itemize}{itemize}{2} % Ograniczenie liczby poziomów do dwóch
\renewlist{enumerate}{enumerate}{2} % Ograniczenie liczby poziomów do dwóch
%------------------------------ Podpisy tabel i ilustracji ------------------------------

\captionsetup{
    width=0.9\textwidth, % Większa szerokość podpisu
    indention=-30pt % Ujemne wcięcie pierwszego wiersza
}

%-------------------------------------- Układ strony --------------------------------------

\fancypagestyle{plain}{
	\fancyhf{}
	\renewcommand{\headrulewidth}{0pt}
	\fancyfoot[LE,RO]{\thepage}
}

\setstretch{1.15}
%------------------------------ Środowiska "abstrakt" ------------------------------
% Streszczenie po polsku
%-----------------------------------------------------------------------------------
\newenvironment{pl_abstract}{
    \thispagestyle{plain} % Styl strony
    \par % Rozpocznij paragraf
    \large % Ustaw czcionkę
    \vspace*{1cm} % Odstęp pod tytułem dokumentu
    \begin{center}
        \textbf{\Large
        Niepotrzebnie długi i skomplikowany tytuł pracy \\ trudny do przeczytania, zrozumienia i wymówienia
        } % Wycentrowany pogrubiony tytuł dokumentu
    \end{center}
    \vspace*{4cm} % Odstęp pod tytułem dokumentu
    \begin{center}
        \noindent\textbf{\large Abstrakt}
    \end{center}
    % \vspace{0.5cm} % Dodatkowy odstęp pod "Streszczenie"
    \par\noindent\ignorespaces
}{
    \ignorespacesafterend % Usuń odstępy po środowisku
}
\newenvironment{slowa_klucze}{
    % \par % Rozpoczyna nowy akapit
    \large % Lokalna zmiana czcionki
    \thispagestyle{plain} % Styl strony
    \vspace{0.5cm} % Odstęp
    \noindent \textbf{Słowa kluczowe:} 
    \ignorespaces % Ignoruje białe znaki przed treścią środowiska
}{
    \par % Kończy akapit
    \unskip % Usuwa białe znaki na końcu
}
%-----------------------------------------------------------------------------------
% Streszczenie po angielsku
%-----------------------------------------------------------------------------------
\newenvironment{eng_abstract}{
    \thispagestyle{plain} % Styl strony
    \par % Rozpocznij paragraf
    \large % Ustaw czcionkę
    \vspace*{1cm} % Odstęp pod tytułem dokumentu
    \begin{center}
        \textbf{\Large
        Unnecessarily long and complicated thesis' title \\ difficult to read, understand and pronounce
        } % Wycentrowany pogrubiony tytuł dokumentu
    \end{center}
    \vspace*{4cm} % Odstęp pod tytułem dokumentu
    \begin{center}
        \noindent\textbf{\large Abstract}
    \end{center}
    \par\noindent\ignorespaces
}{
    \ignorespacesafterend % Usuń odstępy po środowisku
}
\newenvironment{second_keywords}{
    \par % Rozpoczyna nowy akapit
    \large % Lokalna zmiana czcionki
    \thispagestyle{plain} % Styl strony
    \vspace{0.5cm} % Odstęp
    \noindent \textbf{Keywords:} 
    \ignorespaces % Ignoruje białe znaki przed treścią środowiska
}{
    \par % Kończy akapit
    \unskip % Usuwa białe znaki na końcu
}


%------------------------------ Spis załączników ------------------------------
\makeatletter
% Przedefiniowanie zachowania \appendix
\newcommand{\nocontentsline}[3]{}% Wyłączenie linii dodawania do spisu treści
\newcommand{\tocless}[2]{% Tworzenie załącznika bez dodawania do spisu treści
    \bgroup\let\addcontentsline=\nocontentsline#1{#2}\egroup
}
\newcommand{\@appendixtitle}{% Nagłówek dla załącznika
        Załącznik
}
\let\standardappendix\appendix
\renewcommand\appendix[1]{
    \refstepcounter{section}
    \renewcommand\thesection{\Alph{section}} % Numeracja literowa
    \tocless{\section*}{\@appendixtitle~\thesection\hspace*{2.5em} #1}
    \addcontentsline{app}{subsection}{\hspace*{0em}\thesection.\hspace*{0.5em} #1}
}
\newcommand\listofappendicestoc{
    \section*{Spis załączników}\@starttoc{app}
    \standardappendix
    \renewcommand{\thesection}{\arabic{section}}
}
\newcommand{\listoffigurestoc}{
    \listoffigures
    \addcontentsline{toc}{section}{\listfigurename}
}
\newcommand{\listoftablestoc}{
    \listoftables
    \addcontentsline{toc}{section}{\listtablename}
}

%------------------------------ Metadane dokumentu ------------------------------
\hypersetup{
    colorlinks,
    citecolor=black,
    filecolor=black,
    linkcolor=black,
    urlcolor=black,
    pdftitle={Temat pracy dyplomowej},
    pdfsubject={Zagadnienia poruszane w pracy},
    pdfauthor={Autor},
    pdfcreator={TeX/LaTeX},
    pdfproducer={MikTeX wersja ..., pdflatex wersja ...},
    pdfkeywords={Słowa kluczowe}}

\graphicspath{{img/}} % Katalog z ilustracjami           

\usepackage{strona_tytulowa_pliki/FrontPage}

%----------------------------------- Akronimy -----------------------------------

\DeclareAcronym{pw}{short=PW, long=Politechnika Warszawska}
\DeclareAcronym{ei}{short=EiTI, long=Wydział Elektroniki i Technik Informacyjnych}

%----------------------------------- Częściowa kompilacja -----------------------

% Odkomentowanie poniższej linii spowoduje kompilację 
% tylko wybranych plików. Stare odwołania i 
% numeracja zostaną zachowane. Przydatne w przypadku dużych prac.
% \includeonly{tex/0-abstrakt}          % Częściowa kompilacja 
% \includeonly{tex/1-czesc-teoretyczna} % Częściowa kompilacja 
% \includeonly{tex/2-czesc-wykonawcza}  % Częściowa kompilacja 
% \includeonly{tex/3-podsumowanie}      % Częściowa kompilacja 

\begin{document}
%-------------------------Strona tytułowa------------------------------

% Poniższe pola należy uzupełnić odpowiednimi danymi
\instytut{\{nazwa jednostki\}}
\typPracy{i} % i dla pracy inżynierskiej, m dla pracy magisterskiej
\kierunek{\{nazwa kierunku\}}
\specjalnosc{\{nazwa specjalności\}}
\title{\{Tytuł pracy\}}
\author{\{Imię i Nazwisko\}}
\album{\{liczba\}}
% Istnieje możliwość dodania drugiego autora. W tym celu należy użyć poniższej funkcji. Nie można dodać imienia autora bez podania jego numeru albumu.
% \DrugiAutor{\{Imię i nazwisko\}}
% \DrugiAlbum{\{liczba\}}
\promotor{\{tytuł/stopień naukowy, Imię i Nazwisko\}}
% Istnieje możliwość dodania opiekuna oraz konsultanta
% \opiekun{\{tytuł/stopień naukowy, Imię i Nazwisko\}}
% \konsultacje{\{tytuł/stopień naukowy, Imię i Nazwisko\}}
\date{\the\year}
\maketitle{}
%------------------------------------------------

\cleardoublepage % Zaczynamy od nieparzystej strony
\begin{pl_abstract}
    As any dedicated reader can clearly see, the Ideal of practical reason is a re-
    presentation of, as far as I know, the things in themselves; as I have shown elsewhere, the
    phenomena should only be used as a canon for our understanding. The paralogisms of
    practical reason are what first give rise to the architectonic of practical reason. As will easily
    be shown in the next section, reason would thereby be made to contradict, in view of these
    considerations, the Ideal of practical reason, yet the manifold depends on the phenomena.
    Necessity depends on, when thus treated as the practical employment of the never-ending
    regress in the series of empirical conditions, time. Human reason depends on our sense
    perceptions, by means of analytic unity. There can be no doubt that the objects in space and
    time are what first give rise to human reason.
    Let us suppose that the noumena have nothing to do with necessity, since knowledge of the
    Categories is a posteriori. Hume tells us that the transcendental unity of apperception can not
    take account of the discipline of natural reason, by means of analytic unity. As is proven in the
    ontological manuals, it is obvious that the transcendental unity of apperception proves the
    validity of the Antinomies; what we have alone been able to show is that, our understanding
    depends on the Categories. It remains a mystery why the Ideal stands in need of reason. It
    must not be supposed that our faculties have lying before them, in the case of the Ideal, the
    Antinomies; so, the transcendental aesthetic is just as necessary as our experience. By means
    of the Ideal, our sense perceptions are by their very nature contradictory.
\end{pl_abstract}

\begin{slowa_klucze}
    jakieś, słowa, klucze, jakieś, słowa, klucze, jakieś, słowa, klucze, jakieś, słowa, klucze, jakieś, słowa, klucze,
    jakieś, słowa, klucze, jakieś, słowa, klucze, jakieś, słowa, klucze, jakieś, słowa, klucze, jakieś, słowa, klucze,
    jakieś, słowa, klucze, jakieś, słowa, klucze, jakieś.
\end{slowa_klucze}

%----------------------------------------
% Streszczenie po angielsku
%----------------------------------------
\clearpage
\begin{eng_abstract}
    As any dedicated reader can clearly see, the Ideal of practical reason is a re-
    presentation of, as far as I know, the things in themselves; as I have shown elsewhere, the
    phenomena should only be used as a canon for our understanding. The paralogisms of
    practical reason are what first give rise to the architectonic of practical reason. As will easily
    be shown in the next section, reason would thereby be made to contradict, in view of these
    considerations, the Ideal of practical reason, yet the manifold depends on the phenomena.
    Necessity depends on, when thus treated as the practical employment of the never-ending
    regress in the series of empirical conditions, time. Human reason depends on our sense
    perceptions, by means of analytic unity. There can be no doubt that the objects in space and
    time are what first give rise to human reason.
    Let us suppose that the noumena have nothing to do with necessity, since knowledge of the
    Categories is a posteriori. Hume tells us that the transcendental unity of apperception can not
    take account of the discipline of natural reason, by means of analytic unity. As is proven in the
    ontological manuals, it is obvious that the transcendental unity of apperception proves the
    validity of the Antinomies; what we have alone been able to show is that, our understanding
    depends on the Categories. It remains a mystery why the Ideal stands in need of reason. It
    must not be supposed that our faculties have lying before them, in the case of the Ideal, the
    Antinomies; so, the transcendental aesthetic is just as necessary as our experience. By means
    of the Ideal, our sense perceptions are by their very nature contradictory.
\end{eng_abstract}

\begin{second_keywords}
    some, words, keywords, some, words, keywords, some, words, keywords, some, words, keywords, some, words, keywords, 
    some, words, keywords, some, words, keywords, some, words, keywords, some, words, keywords, some, words, keywords, 
    some, words, keywords, some, words, keywords, some, words, keywords.
\end{second_keywords}

\pagestyle{plain}

\cleardoublepage % Zaczynamy od nieparzystej strony
\tableofcontents % Spis treści

% %------------------------------------------------
% % Części pracy
% %------------------------------------------------
\cleardoublepage % Powinna zaczynać się na nieparzystej stronie
\section*{Wprowadzenie}
    \noindent\lipsum[1-3]

    \subsection*{Cel pracy}
        \noindent\lipsum[4-5]

\cleardoublepage
\part{Analiza stanu wiedzy i dobór narzędzi}{I}

\noindent\lipsum[1-3]
\clearpage\section{Istniejące rozwiązania}
\noindent\lipsum[1-3]

 % Przykładowy obrazek
\begin{figure}[!h]
    % Wyrównanie obrazka, szerokość i plik
    % Zamiast width można też użyć height, etc.
    \centering \includegraphics[width=0.5\linewidth]{logopw.png}
    % Podpis umieszczamy pod obrazkiem
    % znacznik \caption służy również do wygenerowania numeru obrazka
    
    \caption{Lorem ipsum dolor sit amet, consectetuer adipiscing elit. Ut purus
    elit, vestibulum ut, placerat ac, adipiscing vitae, felis}
    % \label pozwala odwołać się do obrazka w innych miejscach za pomocą \ref
    % odwołanie \ref renderuje się jako numer obrazka,
    % dlatego zawsze najpierw używaj \caption a potem \label
    \label{fig:tradycyjne-logo-pw}
\end{figure}

% \ref wyrenderuje się jako 'Reference to image 1.1'
\lipsum[1-3] Reference to image \ref{fig:tradycyjne-logo-pw}.

% Przypis dolny \footnote
\lipsum[4] Lorem ipsum dolor sit amet\footnote{Lorem ipsum dolor sit amet, consectetur adipiscing elit, sed do eiusmod tempor incididunt ut labore et dolore magna aliqua. Ut enim ad minim veniam, quis nostrud exercitation ullamco laboris nisi ut aliquip ex ea commodo consequat.}, consectetur adipiscing elit.

\subsection{Przegląd urządzeń o podobnym zastosowaniu}

\noindent Poniżej znajdują się przykładowe odwołania do bibliografii. 

%Ta lista jest długa. Zmnieszamy czcionkę i odstępy.
% \begingroup

\begin{itemize} \footnotesize
    \item książka (jako całość) z minimum wymaganych pól \cite{Gorski2023},
    \item książka (jako całość) z kilkoma autorami \cite{Kurczab2016},
    \item książka (jako całość) z informacją o wydaniu oraz ISBN \cite{Rowling2000},
    \item książka (jako całość) z informacją o tłumaczu. \cite{AlKhalili2022},
    \item rodzdział książki bez wydzielonego autorstwa \cite{Smith2018},
    \item rodzdział książki bez wydzielonego autorstwa \cite{Peterson2019},
    \item rozdział książki z autorami rozdziału \cite{Kraus2010},
    \item rozdział książki z autorami rozdziału \cite{Isidori2023},
    \item rozdział książki z autorami rozdziału \cite{Gregerson2010},
    \item artykuł w czasopiśmie \cite{Nowak2023},
    \item artykuł w czasopiśmie \cite{Xu2023},
    \item artykuł w czasopiśmie \cite{Mashatan2007},
    \item artykuł w czasopiśmie \cite{Hwang2006},
    \item artykuł w czasopiśmie \cite{Joye2007},
    \item praca dyplomowa \cite{Cwik2018},
    \item praca dyplomowa \cite{Marszalek2023},
    \item praca dyplomowa \cite{Urbanski2023},
    \item praca dyplomowa \cite{Otto2004},
    \item dokumentacja techniczna \cite{LPC2212},
    \item dokumentacja techniczna \cite{NIST-SP-800-22},
    \item standard \cite{FIPS140-2},
    \item akt prawny \cite{ust2001-09-18},
    \item publikacja internetowa \cite{Schneier2017},
    \item publikacja internetowa \cite{Craddock2017},
    \item publikacja internetowa \cite{MatlabPerformance2015},
    \item materiały konferencyjne \cite{Duckworth2019},
    \item materiały konferencyjne \cite{Jun2006}.
\end{itemize}
% \endgroup

\subsection{Wady i zalety obecnych rozwiązań}
\noindent Poniżej znajdują się przykłady list. Listy punktowe używane są do wymieniania elementów, 
nieocechowanych priorytetem. Według zaleceń Rektora Politechniki Warszawskiej, w pracach
dyplomowych zalecane jest, aby wyliczenia punktowe symbolizowane były kropką
lub myślnikiem (--).

\noindent Listy punktowe -- (przykład z „kropką”):
\begin{itemize}
    \item pierwszy punkt,
    \item drugi punkt,
    \item trzeci punkt.
\end{itemize}

\noindent Listy punktowe. Przykład z „--” oraz zagnieżdżeniem: 
\begin{itemize}[label=---]
    \item[--] Pierwszy punkt z myślnikiem.
    \item[--] Drugi punkt z myślnikiem.
    \item[--] Trzeci punkt z myślnikiem.
    \begin{itemize}
        \item[--] Drugi pierwszy punkt z myślnikiem.
        \item[--] Drugi drugi punkt z myślnikiem.
        \item[--] Drugi trzeci punkt z myślnikiem.
        \begin{itemize}
            \item[--] Pierwszy drugi pierwszy punkt z myślnikiem.
            \item[--] Drugi drugi drugi punkt z myślnikiem.
            \item[--] Trzeci drugi trzeci punkt z myślnikiem.
        \end{itemize}    
    \end{itemize} 
    
    \item[--] Trzeci punkt z myślnikiem.
\end{itemize}

% Lista punktowana
% Parametr label ustawia symbol punktora
\noindent Należy mieć na uwadze, że istnieje wiele różych „kresek”, które wyglądają jak 
myślnik, pachną jak myślnik, ale myślnikiem nie są. Oto one:
\begin{itemize}[label=--]
    \item dywiz „-”
    \item myślnik lub półpauza „--”,
    \item pauza „---”
    \item znak minus „$-$”
\end{itemize}   
Każdy z tych znaków ma swoje unikalne przeznaczenie. 

\begin{table}[H] \centering
    \caption{Przykładowa tabela w stylu IEEE}
    \label{tab:IEEE}
    \begin{tabular}{c c c c c}
        \toprule[2pt]
         & \multicolumn{2}{c}{\textbf{\uline{Cell}}}& \multicolumn{2}{c}{\textbf{\uline{Cell}}}   \\
        \textbf{{Cell}}  &\textbf{{Cell}}     &  \textbf{Cell}&   \textbf{Cell} &  \textbf{Cell} \\  \midrule
        cell  & cell & cell & cell & cell   \\
        cell  & cell & cell & cell & cell   \\
        cell  & cell & cell & cell & cell   \\
        cell  & cell & cell & cell & cell   \\
        cell  & cell & cell & cell & cell   \\ 
        cell  & cell & cell & cell & cell   \\ \midrule
        \multicolumn{4}{r}{\textbf{Cell}} & Cell \\
        \bottomrule[2pt]
    \end{tabular}
\end{table}

\noindent Listy wyliczeniowe. Stosowane są w przypadku uporządkowanych lub upriorytetowanych elementów. Zazwyczaj symbolizujemy je liczbami arabskimi, rzymskimi czy też literami.
\begin{enumerate}
    \item Punkt pierwszy.
    \item Punkt drugi.
    \item Punkt trzeci.
\end{enumerate}
\noindent Listy definicyjne:
Używane są do objaśniania wszelkiego rodzaju pojęć,  akronimów czy też symboli, występują w dwóch
wariacjach -- z wcięciami w formie zdaniowej lub z wyrównaniem kolumnowym (w postaci tabeli).
W tym szablonie jest jednak odrębne miejsce na umieszczenie własnych akronimów i symboli. Znajduje
się ono w pliku main.tex.

\begin{description}
    \item[Kot] ma 4 nogi, miauczy i bardzo dużo śpi w ciągu dnia.
    \item[Pies] to czteronożny zwierz, który szczeka, towarzyszy człowiekowi w 
   spacerach oraz aportuje piłki.
    \item[Świnia] ma 4 nogi, kwiczy i jest zwykle zwierzęciem chlewnym, które jest 
   czasem wypowadzane na spacery.
\end{description}
\noindent Ogólne zasady tworzenia list:
\begin{itemize}
    \item Każda lista powinna zawierać co najmniej 3 elementy.
    \item Należy unikać zbyt długich wyliczeń lub wypunktowań, w takim wypadku rozważyć podzielenie 
    na mniejsze listy lub podlisty.
    \item Elementy wyliczeń, które nie przybierają postaci samodzielnych zdań rozpoczynamy małą literą.
    \item Całą listę, tj. ostatni element wyliczania, zawsze kończymy kropką.
\end{itemize}


\clearpage\section{Normy i standardy techniczne}
\noindent \lipsum[1-3]

    \subsection{Wymagania prawne}
    \noindent \lipsum[4-5]

    \subsection{Standardy projektowe}
    \noindent \lipsum[6-7]

\clearpage\section{Określenie problemu}
    \noindent \lipsum[8-9]

\clearpage\section{Dobór narzędzi}
    \noindent \lipsum[1-3]

    \subsection{Projekt 3D}
        \noindent \lipsum[4-5]

    \subsection{Schematy elektryczne}
        \noindent \lipsum[6-7]
        
    \subsection{Środowisko symulacyjne}
        \noindent \lipsum[8-9]


\clearpage\section{Wymagania funkcjonalne i techniczne}
    \noindent \lipsum[1-3]

    \subsection{Specyfikacja funkcjonalna}
        \noindent \lipsum[4-5]

    \subsection{Ograniczenia techniczne}
        \noindent \lipsum[6-7]
\cleardoublepage % Powinna zaczynać się na nieparzystej stronie
\part{Projekt i implementacja}{II}
\noindent \lipsum[1-3]

\clearpage\section{Koncepcja urządzenia}
    % Równanie typu 'inline':
    \noindent \lipsum[1-3] $F = m \cdot a$ lorem ipsum dolor sit amet.
    % Równanie bez numeru
    % align oznacza wyrównanie kolejnych wierszy do '&'
    % '&' służy tylko do wyrównania i nie jest renderowany
    \begin{align}
        E & = mc^2 \\
        y & = ax^2 + bx + c
    \end{align}

    \lipsum[4-5]
    % Równanie numerowane: macierze
    \begin{align}
        \begin{bmatrix}
            1 & 0 & 0 \\
            0 & 2 & 0 \\
            0 & 0 & 3
        \end{bmatrix} \cdot
        \begin{bmatrix}
            4 \\
            5 \\
            6
        \end{bmatrix} =
        \begin{bmatrix}
            4  \\
            10 \\
            18
        \end{bmatrix}
    \end{align}


    \subsection{Szkic koncepcyjny}
        \noindent \lipsum[6-7]

    \subsection{Schemat blokowy}
        \noindent \lipsum[8-9]

\clearpage\section{Dobór technologii i komponentów} 
    \noindent \lipsum[1-3]

    \subsection{Komponenty mechaniczne}
        \noindent \lipsum[4-5]

    \subsubsection{Układy elektroniczne}
        \noindent \lipsum[6-7]

    \subsubsection{Oprogramowanie sterujące} 
        \noindent \lipsum[8-9]

\clearpage\section{Symulacje i analizy projektowe}
    \subsection{Analiza wytrzymałości konstrukcji}    
        \noindent \lipsum[1-3]
    
    \subsection{Symulacja układu elektronicznego}    
        \noindent \lipsum[4-5]
    
\clearpage\section{Etapy realizacji}
    \noindent \lipsum[1-3]
    \subsection{Montaż mechaniczny}
        \noindent \lipsum[4-5]

    \subsection{Montaż elektroniczny}
        \noindent \lipsum[6-7]

    \subsection{Integracja systemu}
        \noindent \lipsum[8-9]

\clearpage\section{Dokumentacja techniczna}
    \noindent \lipsum[1-5]
\cleardoublepage % Powinna zaczynać się na nieparzystej stronie
\part{Badania i analiza powykonawcza}{III}
\noindent \lipsum[1-5]

\clearpage\section{Metodyka testowania}
\noindent \lipsum[6-8]

% Fragment kodu źródłowego programu
% \addmargin pozwala na wcięcie kodu od lewej (tu: 8mm).
% Wcięcie służy do tego, aby numery linii nie wystawały poza lewy margines.
% Druga liczba oznacza wcięcie od prawej.
\begin{addmargin}[8mm]{0mm}
\begin{lstlisting}[
    language=HTML,
    numbers=left,
    firstnumber=1,
    caption={\textbf{Hello world} w HTML},
    aboveskip=10pt
]
<html>
  <head>
    <title>Hello world!</title>
  </head>
  <body>
    Hello world!
  </body>
</html>
\end{lstlisting}
\end{addmargin}

\lipsum[9-10]

% Dla dłuższych numerów linii potrzebne jest większe wcięcie.
\begin{addmargin}[12mm]{0mm}
\begin{lstlisting}[
    language=C++,
    numbers=left,
    firstnumber=147,
    caption={Generowanie sekwencji Collatza w języku C++},
    aboveskip=10pt
]
class Collatz {
  private:
    unsigned current_val_;
    void update_val() {
        if( current_val_ % 2 == 0 )
            current_val_ /= 2;
        else
            current_val_ = current_val_ * 3 + 1;
    }

  public:
    explicit Collatz(unsigned initial_value) :
        current_val_(initial_value) {}
    void print_sequence() {
        unsigned i = 1;
        while( current_val_ > 1 ) {
            std::cout
                << "val " << i << " = " << current_val_
                << std::endl;
            update_val(); ++i;
        }
    }
};

int main() {
  // prints Collatz seqence, starting from 194375
  Collatz seq(194375);
  seq.print_sequence();
  return 0;
}
\end{lstlisting}
\end{addmargin}

\lipsum[10-12]
 
    \subsection{Testy funkcjonalne}
      \noindent\lipsum[1-3]

    \subsection{Testy wytrzymałościowe}
      \noindent\lipsum[4-7]

    \subsection{Testy wydajnościowe}
      \noindent\lipsum[8-10]

\clearpage\section{Wyniki testów}
  \noindent\lipsum[1-7]

    \subsection{Porównanie wyników z założeniami}
      \noindent\lipsum[1-3]

    \subsection{Analiza błędów}
      \noindent\lipsum[4-6]

\clearpage\section{Analiza ekonomiczna}  
  \noindent\lipsum[1-7]

  \subsection{Koszty materiałów i komponentów}
    \noindent\lipsum[1-3]

  \subsection{Koszty produkcji seryjnej}
    \noindent\lipsum[4-7]

  \subsection{Porównanie kosztów z alternatywami}
    \noindent\lipsum[8-10]

\clearpage\section{Problemy napotkane podczas budowy i ich rozwiązanie}
\noindent \lipsum[1-5]

\clearpage\section*{Podsumowanie}
    \noindent \lipsum[1-9]
%------------------------------------------------
\cleardoublepage % Powinna zaczynać się na nieparzystej stronie
\printbibliography % Bibliografia
\clearpage

% Wykaz symboli i skrótów.
% Należy posortować symbole alfabetycznie
% we własnym zakresie.
\printacronyms % Wydruk listy akronimów (opcjonalne)


%--------------------------------------
% Spisy
%--------------------------------------
\pagestyle{plain}
\newpage   
\listoffigurestoc       % Spis rysunków
\newpage                % Odstęp pionowy
\listoftablestoc        % Spis tabel
\newpage                % Odstęp pionowy
\listofappendicestoc    % Spis załączników

% %--------------------------------------
% % Załączniki
% %--------------------------------------
% % Ilustracje i tabele w załącznikach nie pojawią się
% % w głównym spisie treści
\captionsetup[figure]{list=no}
\captionsetup[table]{list=no}
% Załącznik 1

\clearpage
\appendix{Nazwa załącznika A}
\noindent\lipsum[1-3]
\begin{figure}[!h]
	\centering \includegraphics[width=0.5\linewidth]{logopw2.png}
	\caption{Obrazek w załączniku}
\end{figure}
\lipsum[4-7]

% Załącznik 2
\clearpage
\appendix{Nazwa załącznika B}
\lipsum[1-2]
\begin{table}[H] \centering
    \caption{Opis wyprowadzeń i sygnałów}
    \label{tab:base_pins}
    \begin{tabular}{c c c c c}
        \toprule[2pt]
         & \multicolumn{2}{c}{\textbf{\uline{Cell}}}& \multicolumn{2}{c}{\textbf{\uline{Cell}}}   \\
        \textbf{{Cell}}  &\textbf{{Cell}}     &  \textbf{Cell}&   \textbf{Cell} &  \textbf{Cell} \\  \midrule
        cell  & cell & cell & cell & cell   \\
        cell  & cell & cell & cell & cell   \\
        cell  & cell & cell & cell & cell   \\
        cell  & cell & cell & cell & cell   \\
        cell  & cell & cell & cell & cell   \\ 
        cell  & cell & cell & cell & cell   \\ \midrule
        \multicolumn{4}{r}{\textbf{Cell}} & Cell \\
        \bottomrule[2pt]
    \end{tabular}
\end{table}
\lipsum[3-4]



\end{document}